\documentclass[11pt]{article}
\usepackage[utf8]{inputenc}
\usepackage[T1]{fontenc}
\usepackage[spanish]{babel}
\usepackage[margin=2.2cm]{geometry}
\usepackage{booktabs}
\usepackage{array}
\usepackage{amsmath}
\usepackage{amssymb}
\usepackage{xcolor}
\usepackage[colorlinks=true,linkcolor=blue,citecolor=blue,urlcolor=blue]{hyperref}
\usepackage{enumitem}
\usepackage{titlesec}
\usepackage{parskip}

\titleformat{\section}{\normalfont\Large\bfseries}{\thesection.}{0.6em}{}
\titlespacing*{\section}{0pt}{1.4em}{0.6em}

\title{\textbf{El Clasificador}\\[0.2em]\large Aprendizaje por Refuerzo para la detección de correo anómalo}
\author{Sistemas Inteligentes \\ Universidad de La Frontera}
\date{}

\begin{document}
\maketitle
\vspace{-2em}
\begin{center}
\small Informe breve \quad\textbullet\quad Aplicación web: \textit{El Clasificador}
\end{center}
\vspace{1em}

\noindent\textbf{Resumen.} Este trabajo presenta una aplicación web educativa
que enseña Aprendizaje por Refuerzo (RL) a través de un caso de detección de
anomalías: separar correo legítimo de correo anómalo (spam y suplantación,
\textit{phishing}). En lugar de un tutorial abstracto, la app muestra a un
agente de Q-learning (un clasificador recién instalado en una sala de correo)
fracasar, aprender y finalmente superar a un conjunto de reglas razonables.
Cada concepto formal (estado, acción, recompensa, política) aparece primero
como un hecho físico y solo después como una ecuación. Todo el modelo se
entrena en el navegador, sin backend ni dependencias.

\section{El problema, y por qué es Aprendizaje por Refuerzo}

Clasificar correo es tomar, miles de veces al día, una decisión bajo
incertidumbre: \texttt{aceptar} (al buzón) o \texttt{descartar} (a la
papelera), sin conocer con certeza el contenido. Los errores no son
simétricos: dejar pasar spam es una molestia. \textbf{Descartar una carta
legítima es un daño irreversible.} Esa asimetría de costos, sumada a la
naturaleza secuencial del flujo (la reputación de un remitente evoluciona con
lo que ha enviado antes), es exactamente la clase de problema que el RL
modela de forma natural: un agente que aprende una política por interacción
con un entorno, guiado solo por una señal de recompensa \cite{sutton2018}.

A diferencia del aprendizaje supervisado, aquí no hay un maestro que entregue
la etiqueta correcta de cada correo antes de decidir. Hay un número, la
recompensa, que llega \textit{después} de actuar, cuando se revela la verdad,
y que empuja al agente a mejorar. El costo asimétrico, además, no hay que
codificarlo como una regla: vive dentro de la función de recompensa, y la
política conservadora emerge.

\textbf{¿Por qué esto es detección de anomalías?} La detección de anomalías
identifica eventos poco frecuentes que se desvían de un patrón esperado,
mayormente ``normal'' \cite{pang2021}. El correo anómalo de este trabajo
encaja ahí: spam y suplantación son minoría (30\% y 10\% del tráfico, ver
\S4) frente al correo legítimo, y se definen precisamente por desviarse de
las señales típicas de un remitente de confianza: contacto nuevo o
reputación en caída, urgencia inusual, enlaces sin patrón previo, horario
atípico. La suplantación es el caso límite: no es una categoría aparte con
reglas propias, es la \textit{cola} de la distribución, camuflada dentro de
lo que parece normal (buena reputación, remitente conocido). Detectarla
exige aprender la combinación de señales, no una sola regla. Es exactamente
lo que el agente hace.

\section{El modelo: un proceso de decisión (MDP)}

Formalizamos la sala de correo como una tupla $(S, A, R, T, \gamma)$.

\textbf{Estado ($S$).} Cada correo se describe con cinco variables
observables, dibujadas sobre el sobre:

\begin{center}
\begin{tabular}{@{}lll@{}}
\toprule
\textbf{Variable} & \textbf{Valores} & \textbf{Representación visual} \\
\midrule
remitente & conocido / desconocido & sello relleno o hueco \\
reputación & buena / neutra / mala & color del sello (ámbar / gris / negro) \\
enlaces & 0 / 1 / varios & marcas de anzuelo \\
urgencia & baja / alta & cinta roja \\
hora & laboral / madrugada & matasellos (sol / luna) \\
\bottomrule
\end{tabular}
\end{center}

Cardinalidad: $2 \times 3 \times 3 \times 2 \times 2 = 72$ estados. La
discretización se eligió para ser legible, no óptima: permite representar
toda la tabla Q como un archivador de 72 sobres que el usuario puede leer de
un vistazo.

\textbf{Acciones ($A$).} $\{\text{aceptar}, \text{descartar}\}$.

\textbf{Recompensa ($R$).} Es literalmente la matriz de confusión, con costo
asimétrico:

\begin{center}
\begin{tabular}{@{}lrr@{}}
\toprule
 & \textbf{era legítimo} & \textbf{era spam} \\
\midrule
\textbf{aceptar} & $+5$ & $-2$ \\
\textbf{descartar} & $\mathbf{-20}$ & $+3$ \\
\bottomrule
\end{tabular}
\end{center}

Descartar una carta real cuesta diez veces más que dejar pasar spam. Ese
$-20$ es el corazón del modelo.

\textbf{Transición ($T$).} Tras cada decisión se revela la etiqueta
verdadera y el entorno actualiza la reputación del remitente: sube un nivel
con correo legítimo (lentamente: dos aciertos por nivel), baja con spam (un
envío por nivel). Así el estado de un remitente evoluciona con la
trayectoria.

\textbf{Descuento ($\gamma = 0.9$).}

\section{El algoritmo: Q-learning tabular}

El agente estima $Q(s, a)$, el valor a largo plazo de tomar la acción $a$ en
el estado $s$, y lo corrige tras cada transición $(s, a, r, s')$
\cite{sutton2018,watkins1992}:

\begin{center}
$Q(s,a) \leftarrow Q(s,a) + \alpha \cdot \left[\, r + \gamma \cdot
\max_{a'} Q(s',a') - Q(s,a) \,\right]$
\end{center}

La actualización es \textit{off-policy}: el objetivo usa
$\max_{a'} Q(s',a')$, la mejor acción futura, aunque el agente aún esté
explorando. Hiperparámetros: $\alpha = 0.1$, $\gamma = 0.9$, y una política
$\varepsilon$-greedy con $\varepsilon: 1.0 \rightarrow 0.05$ (decaimiento
exponencial sobre 5\,000 episodios). No se usa una red neuronal (DQN): el
estado es diminuto, y una tabla, a diferencia de una red, se puede
visualizar, que es el punto de la app.

\section{Los datos: un entorno con trampas}

El correo es sintético y se muestrea de cuatro poblaciones con
características estocásticas (por eso ninguna señal aislada clasifica
perfecto): legítimo cotidiano (45\%), legítimo nuevo (15\%), spam burdo
(30\%) y suplantación (10\%). Un episodio es una jornada de 50 sobres de 15
remitentes fijos. El entorno tiene tres trampas deliberadas:

\begin{itemize}[leftmargin=1.4em]
  \item \textbf{El legítimo nuevo} castiga la regla ``descartar si el
    remitente es desconocido'': rompe cartas reales de clientes nuevos.
  \item \textbf{La suplantación} castiga la regla ``confiar en el remitente
    conocido'': llega desde una dirección conocida y de buena reputación.
    Solo la delatan los enlaces.
  \item \textbf{El remitente comprometido} (3 de los 15) empieza enviando
    correo legítimo y, a mitad de la jornada, pasa a la suplantación. Impide
    memorizar ``confío en X''.
\end{itemize}

\section{Resultados}

Se evaluaron cinco estrategias sobre 500 episodios (25\,000 sobres) del
mismo flujo (el entorno es determinista dado su semilla, y su flujo no
depende de la acción, lo que hace la comparación exacta):

\begin{center}
\begin{tabular}{@{}lrrrr@{}}
\toprule
\textbf{Estrategia} & \textbf{R/ep} & \textbf{Cartas rotas} & \textbf{Suplant. atrapada} & \textbf{Aciertos} \\
\midrule
Aceptar todo & 108.9 & 0 & 0\% & 60\% \\
Descartar todo & $-536.2$ & 14\,918 & 100\% & 40\% \\
Regla: descartar si desconocido & $-114.6$ & 4\,927 & 19\% & 49\% \\
Regla: descartar si mala reputación & 176.2 & 0 & 13\% & 87\% \\
\textbf{Política aprendida (RL)} & \textbf{193.0} & 53 & \textbf{40\%} & \textbf{94\%} \\
\bottomrule
\end{tabular}
\end{center}

La política aprendida obtiene la mayor recompensa y atrapa tres veces más
suplantación que la regla de reputación (40\% vs. 13\%), rompiendo muy pocas
cartas reales. La clave: la regla de reputación reacciona (solo descarta
después de que el daño bajó la reputación, y no ve venir al remitente
comprometido, cuya reputación aún es buena al atacar), mientras que la
política aprendida anticipa, desconfiando de la combinación de señales
(sello ámbar más varios anzuelos más urgencia) antes que del historial.

\textbf{La parte honesta.} La política aprendida acepta \textit{más} spam
que la regla dura, porque le enseñamos a temerle al $-20$. No es un defecto:
es la función de recompensa haciendo exactamente lo que le pedimos. La app
lo declara en lugar de esconderlo.

\section{Exploración vs. explotación}

La app permite forzar $\varepsilon = 0$ desde el primer episodio. El agente
resultante explota su primera corazonada de inmediato: su curva de
recompensa arranca alta (no paga el precio de equivocarse al azar) pero se
estanca. Atrapa solo un ${\sim}23\%$ de la suplantación (contra ${\sim}40\%$
del agente que exploró) y deja pasar más spam, porque nunca probó descartar
los sobres ambiguos que parecían legítimos. El agente que explora paga caro
al inicio (su curva se hunde en negativo) y por eso descubre después lo que
la primera corazonada no vio. Es el dilema exploración-explotación, mostrado
en una sola curva.

\section{Limitaciones}

\begin{itemize}[leftmargin=1.4em]
  \item \textbf{Datos sintéticos.} El agente aprendió sobre un mundo
    diseñado por nosotros: proporciones, señales y la recompensa $-20$ son
    elecciones. Ningún resultado aquí prueba algo sobre correo real.
  \item \textbf{¿MDP o bandit contextual?} En este entorno la reputación se
    actualiza según la etiqueta verdadera, no según la acción del agente.
    Estrictamente, la acción no cambia el estado siguiente, por lo que el
    problema se acerca más a un \textit{bandit contextual con contexto no
    estacionario}. El interés secuencial reside en que el estado de un
    remitente evoluciona con la trayectoria (el remitente comprometido).
    Preferimos señalarlo a exagerar el rol de $\gamma$.
  \item \textbf{Alternativas.} Un filtro bayesiano o supervisado, con un
    buen corpus etiquetado, probablemente superaría a RL en exactitud pura
    \cite{sahami1998,cormack2008,pang2021}. Su debilidad es justamente el
    costo asimétrico y el desbalance de clases, que hay que inyectar a
    mano. Este trabajo no afirma que RL sea la mejor solución: muestra cómo
    se piensa el problema en términos de RL.
\end{itemize}

\section{Conclusiones}

\textit{El Clasificador} logra el objetivo del enunciado: explica los
fundamentos del Aprendizaje por Refuerzo (agente, entorno, estado, acción,
recompensa, política), el proceso de aprendizaje por recompensas y el
algoritmo de Q-learning, anclando cada concepto a un caso concreto y
realista de detección de anomalías. La visualización (la compuerta que se
inclina, el archivador de 72 sobres que se colorea, la papelera que
conserva sus cicatrices) convierte conceptos abstractos en objetos
observables. Y, sobre todo, el trabajo es honesto acerca de sus límites, lo
que a nuestro juicio demuestra más comprensión que reportar una cifra de
exactitud sin contexto.

\begin{thebibliography}{9}

\bibitem{sutton2018}
Sutton, R. S., \& Barto, A. G. (2018). \textit{Reinforcement Learning: An
Introduction} (2.\textsuperscript{a} ed.). MIT Press.

\bibitem{watkins1992}
Watkins, C. J. C. H., \& Dayan, P. (1992). Q-learning. \textit{Machine
Learning}, 8(3-4), 279-292.

\bibitem{pang2021}
Pang, G., Shen, C., Cao, L., \& van den Hengel, A. (2021). Deep learning for
anomaly detection: A review. \textit{ACM Computing Surveys}, 54(2).

\bibitem{sahami1998}
Sahami, M., Dumais, S., Heckerman, D., \& Horvitz, E. (1998). A Bayesian
approach to filtering junk e-mail. \textit{AAAI Workshop on Learning for
Text Categorization} (Tech. Report WS-98-05).

\bibitem{cormack2008}
Cormack, G. V. (2008). Email spam filtering: A systematic review.
\textit{Foundations and Trends in Information Retrieval}, 1(4), 335-455.

\end{thebibliography}

\end{document}
